\documentclass[11pt]{article}
\usepackage{mathblatt}
\def\mitzone{1}
\begin{document}
\blattkopf{Lagebeziehungen}{Gesamt}
\verzeichniszeile{\verz{zone}{Kennst du schon (Nr.~1–8)} \verztrenn \verz{e1}{1 Punktprobe und Seitenlage (Nr.~9–11)} \verztrenn \verz{e2}{2 Parameter aus der Lagebedingung (Nr.~12–13)} \verztrenn \verz{e3}{3 Gerade und Ebene (Nr.~14–15)} \verztrenn \verz{e4}{4 Lagebefunde im Sachzusammenhang (Nr.~16–18)} \verztrenn \verz{abhaken}{Das kann ich}}
% Zone „Das kennst du schon“ – aus bank/lagebeziehungen/zone.jsonl (zusammenbau v0.1)
\einheitenkopf[zone][Kennst du schon]{Das kennst du schon}
\setcounter{aufgabe}{0}

%% TODO Titel: Ich-kann-Satz fehlt in der Bank (Platzhalter: Kettenname) – Nr. 1
%% TODO Anweisung: ein Satz über der ersten Teilaufgabe fehlt in der Bank – Nr. 1
\begin{aufgabe}{Ebenengleichungen lesen (Koordinatenform, Normalenvektor ablesen, Parameterform)}
\begin{teile}
\teil $E: 2x + 3y - z = 5$ – ein Normalenvektor? $\vec{n} = $ \leerfeld
\teil $E: 4x - \frac{1}{2}y + 3z = -2$ – ein Normalenvektor? $\vec{n} = $ \leerfeld
\end{teile}
\end{aufgabe}

%% TODO Titel: Ich-kann-Satz fehlt in der Bank (Platzhalter: Kettenname) – Nr. 2
%% TODO Anweisung: ein Satz über der ersten Teilaufgabe fehlt in der Bank – Nr. 2
\begin{aufgabe}{Geradengleichungen lesen und den allgemeinen Geradenpunkt bilden (Stützpunkt plus Parameter mal Richtungsvektor)}
\begin{teile}
\teil $g: \vec{x} = (2 | 1 | 0) + t \cdot (1 | 3 | 2)$ – der Punkt für $t = 1$? \leerfeld
\teil $g: \vec{x} = (1 | -2 | 3) + t \cdot (2 | 1 | -1)$ – der allgemeine Geradenpunkt? \leerfeld
\end{teile}
\end{aufgabe}

%% TODO Titel: Ich-kann-Satz fehlt in der Bank (Platzhalter: Kettenname) – Nr. 3
%% TODO Anweisung: ein Satz über der ersten Teilaufgabe fehlt in der Bank – Nr. 3
\begin{aufgabe}{Skalarprodukt berechnen und Orthogonalität erkennen}
\begin{teile}
\teil $(1 | 2 | 4) \cdot (3 | 0 | 1)$ – wie viel? \leerfeld
\teil $(-3 | 0,5 | 2) \cdot (2 | 4 | -1)$ – wie viel? \leerfeld
\end{teile}
\end{aufgabe}

%% TODO Titel: Ich-kann-Satz fehlt in der Bank (Platzhalter: Kettenname) – Nr. 4
%% TODO Anweisung: ein Satz über der ersten Teilaufgabe fehlt in der Bank – Nr. 4
\begin{aufgabe}{Gleichungen lösen}
\begin{teile}
\teil $3k + 6 = 0$ – $k$? $k = $ \leerfeld
\teil $-4t + 1 = 2t - 2$ – $t$? $t = $ \leerfeld
\end{teile}
\end{aufgabe}

%% TODO Titel: Ich-kann-Satz fehlt in der Bank (Platzhalter: Kettenname) – Nr. 5
%% TODO Anweisung: ein Satz über der ersten Teilaufgabe fehlt in der Bank – Nr. 5
\begin{aufgabe}{Terme mit einem Parameter umformen und Fallunterscheidungen führen}
\begin{teile}
\teil $2(a + 1) - a$ – vereinfacht? \leerfeld
\teil $(a + 2) \cdot t = 4$ mit $a \neq -2$ – $t$? $t = $ \leerfeld
\end{teile}
\end{aufgabe}

%% TODO Titel: Ich-kann-Satz fehlt in der Bank (Platzhalter: Kettenname) – Nr. 6
%% TODO Anweisung: ein Satz über der ersten Teilaufgabe fehlt in der Bank – Nr. 6
\begin{aufgabe}{Punkte im räumlichen Koordinatensystem lesen und Bereiche deuten (Koordinatenschranken, Wandmaße)}
\begin{teile}
\teil Liegt $P(2 | 3 | 1)$ im Quader mit $0 \le x \le 4$, $0 \le y \le 5$ und $0 \le z \le 2$? \\ \kreuz{ja} \\ \kreuz{nein}
\teil Eine Wand liegt in der Ebene $y = 0$ und reicht von $x = 0$ bis $x = 6$ und von $z = 0$ bis $z = 2{,}5$. Liegt $S(4 | 0 | -0{,}5)$ auf der Wand? \\ \kreuz{ja} \\ \kreuz{nein}
\end{teile}
\end{aufgabe}

%% TODO Titel: Ich-kann-Satz fehlt in der Bank (Platzhalter: Kettenname) – Nr. 7
%% TODO Anweisung: ein Satz über der ersten Teilaufgabe fehlt in der Bank – Nr. 7
\begin{aufgabe}{Ebenengleichungen lesen (Koordinatenform, Normalenvektor ablesen, Parameterform) – Fehler finden}
\begin{teile}
\teil Ich finde den Fehler: Zur Ebene $E: 2x - y = 4$ gibt Tim den Normalenvektor $(2 | -1)$ an.
\end{teile}
\end{aufgabe}

%% TODO Titel: Ich-kann-Satz fehlt in der Bank (Platzhalter: Kettenname) – Nr. 8
%% TODO Anweisung: ein Satz über der ersten Teilaufgabe fehlt in der Bank – Nr. 8
\begin{aufgabe}{Ebenengleichungen lesen (Koordinatenform, Normalenvektor ablesen, Parameterform) – selbst rechnen}
\begin{teile}
\teil $E: 5x + z = 3$ – ein Normalenvektor? $\vec{n} = $ \leerfeld
\end{teile}
\end{aufgabe}

\clearpage
% Einheit 1 von 4 – Punktprobe und Seitenlage (bank/lagebeziehungen/e1.jsonl, zusammenbau v0.1)
%% TODO Zeitmarke Einheit 1: Marken-Zeile nicht lesbar
%% TODO Zweigzeile Teil 1: was hier gelernt wird (Satz in Schülersprache aus dem Einheitstitel) – Einheit 1
\einheitenkopf[e1]{Einheit 1 von 4 · Punktprobe und Seitenlage}
\zweigzeile{Abitur GK}
\setcounter{aufgabe}{8}

%% TODO Titel: Ich-kann-Satz fehlt in der Bank (Platzhalter: Kettenname) – Nr. 9
%% TODO Anweisung: ein Satz über der ersten Teilaufgabe fehlt in der Bank – Nr. 9
\begin{aufgabe}{Wer und wogegen?}
\begin{teile}
\teil Gegeben sind ein Punkt $P$ und die Ebene $E: 2x + y - z = 4$. Gefragt ist, ob $P$ in $E$ liegt. Wer gegen wen? \\ \kreuz{Punkt gegen Ebene} \\ \kreuz{Gerade gegen Ebene} \\ \kreuz{Punkt gegen Gerade}
\end{teile}
\end{aufgabe}

%% TODO Titel: Ich-kann-Satz fehlt in der Bank (Platzhalter: Kettenname) – Nr. 10
%% TODO Anweisung: ein Satz über der ersten Teilaufgabe fehlt in der Bank – Nr. 10
\begin{aufgabe}{Punktprobe und Seitenlage}
%% TODO \rechenplatz{n}: Schrittzahl des Grundfalls fehlt in der Bank (nur bei fester Schreibform, 2.3 b)
\begin{teile}
\teil $E: 2x - y + z = 3$. Die Koordinaten von $P$ eingesetzt ergeben links $6$. Wo liegt $P$? \\ \kreuz{in $E$} \\ \kreuz{auf der Seite mit größerem Wert} \\ \kreuz{auf der Seite mit kleinerem Wert}
\teil Liegt $P(3 | 1 | 4)$ in $E: 2x - y + z = 7$?
\teil $E: y - 4z = -6$. Weise nach, dass $S(7 | 2 | 2)$ in $E$ liegt. \hfill (Abitur 2018 LK)
\teil $E: \vec{x} = (2 | 0 | 1) + t \cdot (1 | 1 | 0) + s \cdot (0 | 1 | 2)$. Weise nach, dass $P(4 | 5 | 7)$ in $E$ liegt. \hfill (Abitur 2021 GK)
\teil Eine Lärmschutzwand steht auf der Fahrbahn (xy-Ebene) und liegt in der Ebene $W: 5x - 2y = 0$ (1 LE = 1 m). Weise nach, dass der Punkt $K(2 | 5 | 3)$ in $W$ liegt und dass $W$ senkrecht auf der Fahrbahn steht. \hfill (Abitur 2018 GK)
\teil Ein Quader hat die Kante $AE$ mit $A(5 | 0 | 0)$ und $E(5 | 0 | 3)$; die Ebene $W: x + 2y + 2z = 9$ ist gegeben. Ein Lösungsweg lautet: (1) $P(5 | 0 | r)$ mit $0 \le r \le 3$; (2) $5 + 2 \cdot 0 + 2r = 9$. Formuliere eine passende Aufgabenstellung und führe sie zu Ende. \hfill (Abitur 2024 GK)
\teil Ein Körper ist der Teil des ersten Oktanten zwischen den parallelen Ebenen $L_1: x + 2y + z = 2$ und $L_2: x + 2y + z = 6$. Begründe, dass $T(1 | 0{,}5 | 1)$ im Inneren des Körpers liegt. \hfill (Abitur 2023 GK)
\teil Das Dreieck $ABC$ ist der Teil der Ebene $E: 2x + y + z = 8$ mit nichtnegativen Koordinaten. Aussage: Liegt $P(u | v | w)$ im Innern von $ABC$, dann auch $Q(v | u | w)$. Widerlege die Aussage mit einem Gegenbeispiel. \hfill (Abitur 2026 LK)
\end{teile}
\end{aufgabe}

%% TODO Titel: Ich-kann-Satz fehlt in der Bank (Platzhalter: Kettenname) – Nr. 11
%% TODO Anweisung: ein Satz über der ersten Teilaufgabe fehlt in der Bank – Nr. 11
\begin{aufgabe}{Punktprobe und Seitenlage – Fehler finden, Begründen}
\begin{teile}
\teil Ich finde den Fehler: $E: \vec{x} = (1 | 0 | 2) + t \cdot (1 | 2 | 0) + s \cdot (0 | 1 | 3)$, $P(3 | 5 | 8)$. Lea rechnet: aus x $t = 2$, aus z $s = 2$ – also liegt $P$ in $E$.
\teil Begründe, warum ein Punkt, dessen Koordinaten die Gleichung von $E$ erfüllen, in $E$ liegt.
\end{teile}
\end{aufgabe}

\clearpage
% Einheit 2 von 4 – Parameter aus der Lagebedingung (bank/lagebeziehungen/e2.jsonl, zusammenbau v0.1)
%% TODO Zeitmarke Einheit 2: Marken-Zeile nicht lesbar
%% TODO Zweigzeile Teil 1: was hier gelernt wird (Satz in Schülersprache aus dem Einheitstitel) – Einheit 2
\einheitenkopf[e2]{Einheit 2 von 4 · Parameter aus der Lagebedingung}
\zweigzeile{Abitur LK}
\setcounter{aufgabe}{11}

%% TODO Titel: Ich-kann-Satz fehlt in der Bank (Platzhalter: Kettenname) – Nr. 12
%% TODO Anweisung: ein Satz über der ersten Teilaufgabe fehlt in der Bank – Nr. 12
\begin{aufgabe}{Parameter aus der Lagebedingung}
%% TODO \rechenplatz{n}: Schrittzahl des Grundfalls fehlt in der Bank (nur bei fester Schreibform, 2.3 b)
\begin{teile}
\teil Für welche Zahl $k$ liegt $P(1 | 2 | 0)$ in der Ebene $E_k: kx + 3y + z = 8$? Wer gegen wen? \\ \kreuz{Punkt gegen Ebene} \\ \kreuz{Gerade gegen Ebene} \\ \kreuz{Punkt gegen Gerade}
\teil $E_k: kx + 2y - z = 5$ enthält $P(2 | 1 | 3)$. Bestimme $k$. $k = $ \leerfeld
\teil Das Dreieck mit $A(3 | 0 | 0)$, $B(0 | 6 | 0)$ und $C(0 | 0 | 2)$ liegt in einer Ebene der Form $2x + y + kz = 6$. Bestimme $k$. \hfill (Abitur 2023 GK) \\ $k = $ \leerfeld
\teil Die Ebene $L$ enthält die x-Achse und den Punkt $G(3 | 2 | 5)$; ihre Gleichung hat die Form $r \cdot y + s \cdot z = 0$. Gib passende Werte für $r$ und $s$ an. \hfill (Abitur 2023 GK) \\ $r = $ \leerfeld , $s = $ \leerfeld
\teil Für jede Zahl $a$ ist $g_a: \vec{x} = (3 | 1 | -1) + t \cdot (1 | 2a | a)$. Jede Gerade $g_a$ liegt in der Ebene $E: y - 2z = c$. Bestimme $c$. \hfill (Abitur 2026 LK) \\ $c = $ \leerfeld
\teil Die Ebene $E: 2x + 3y - z = 8$ enthält einen Punkt, dessen drei Koordinaten übereinstimmen. Bestimme ihn. \hfill (Abitur 2019 LK) \\ \leerfeld
\end{teile}
\end{aufgabe}

%% TODO Titel: Ich-kann-Satz fehlt in der Bank (Platzhalter: Kettenname) – Nr. 13
%% TODO Anweisung: ein Satz über der ersten Teilaufgabe fehlt in der Bank – Nr. 13
\begin{aufgabe}{Parameter aus der Lagebedingung – Fehler finden, Begründen}
\begin{teile}
\teil Ich finde den Fehler: $E_k: kx + y + 2z = 8$ soll $P(2 | 4 | 1)$ enthalten. Jana setzt $Q(1 | 2 | 2)$ ein: $k + 2 + 4 = 8$, also $k = 2$.
\teil Begründe, warum das Einsetzen eines Punkts $P$ von $E_k: kx + 2y + z = 4$ eine Gleichung für $k$ liefert.
\end{teile}
\end{aufgabe}

\clearpage
% Einheit 3 von 4 – Gerade und Ebene (bank/lagebeziehungen/e3.jsonl, zusammenbau v0.1)
%% TODO Zeitmarke Einheit 3: Marken-Zeile nicht lesbar
%% TODO Zweigzeile Teil 1: was hier gelernt wird (Satz in Schülersprache aus dem Einheitstitel) – Einheit 3
\einheitenkopf[e3]{Einheit 3 von 4 · Gerade und Ebene}
\zweigzeile{Abitur GK}
\setcounter{aufgabe}{13}

%% TODO Titel: Ich-kann-Satz fehlt in der Bank (Platzhalter: Kettenname) – Nr. 14
%% TODO Anweisung: ein Satz über der ersten Teilaufgabe fehlt in der Bank – Nr. 14
\begin{aufgabe}{Gerade und Ebene}
%% TODO \rechenplatz{n}: Schrittzahl des Grundfalls fehlt in der Bank (nur bei fester Schreibform, 2.3 b)
\begin{teile}
\teil Der allgemeine Punkt von $g$ in die Gleichung von $E$ eingesetzt ergibt $4 + 3t - 3t = 4$. Wie liegt $g$ zu $E$? \\ \kreuz{liegt in der Ebene} \\ \kreuz{echt parallel} \\ \kreuz{genau ein Schnittpunkt}
\teil Weise nach, dass $g: \vec{x} = (1 | 2 | 0) + t \cdot (1 | -1 | 1)$ in $E: x + 2y + z = 5$ liegt.
\teil $g: \vec{x} = (1 | 0 | 2) + t \cdot (2 | 1 | 1)$ und $E: x - 2y + z = 1$. Entscheide mit der Merkregel, wie $g$ zu $E$ liegt.
\teil $P(1 | 3 | 8)$ und $Q_t(5 | t | 6)$ mit reellem $t$. Gibt es ein $t$, für das die Gerade $PQ_t$ parallel zur xy-Ebene verläuft? Begründe. \hfill (Abitur 2024 GK)
\teil Der Punkt $S(4 | 3 | 5)$ wird um die Achse $a: \vec{x} = (0 | 0 | 2) + t \cdot (1 | 2 | 0)$ gedreht. $L: x + 2y = 10$; $H$ ist der Schnittpunkt von $a$ und $L$. Begründe, dass $S$ auf einem Kreis in $L$ mit dem Mittelpunkt $H$ läuft, und bestimme $H$. \hfill (Abitur 2026 LK)
\teil Begründe, dass es unendlich viele Ebenen gibt, die keinen Punkt der Form $(a | a | 0)$ enthalten. \hfill (Abitur 2019 LK)
\teil $E: 2x - z = 1$ und für jede Zahl $a$ die Gerade $g_a: \vec{x} = (1 | 0 | a) + r \cdot (a | 2 | 2)$. Untersuche die Lage von $g_a$ zu $E$ in Abhängigkeit von $a$ und gib gegebenenfalls den Schnittpunkt an. \hfill (Abitur 2024 LK)
\end{teile}
\end{aufgabe}

%% TODO Titel: Ich-kann-Satz fehlt in der Bank (Platzhalter: Kettenname) – Nr. 15
%% TODO Anweisung: ein Satz über der ersten Teilaufgabe fehlt in der Bank – Nr. 15
\begin{aufgabe}{Gerade und Ebene – Fehler finden, Begründen}
\begin{teile}
\teil Ich finde den Fehler: $g: \vec{x} = (1 | 2 | 0) + t \cdot (1 | 0 | 1)$, $E: 2x - y + z = 0$. Nils setzt den Stützpunkt ein, $2 - 2 + 0 = 0$, und schreibt: $g$ liegt in $E$.
\teil Begründe: Fällt beim Einsetzen des allgemeinen Geradenpunkts der Parameter heraus und die Gleichung stimmt, liegt die ganze Gerade in der Ebene.
\end{teile}
\end{aufgabe}

\clearpage
% Einheit 4 von 4 – Lagebefunde im Sachzusammenhang (bank/lagebeziehungen/e4.jsonl, zusammenbau v0.1)
%% TODO Zeitmarke Einheit 4: Marken-Zeile nicht lesbar
%% TODO Zweigzeile Teil 1: was hier gelernt wird (Satz in Schülersprache aus dem Einheitstitel) – Einheit 4
\einheitenkopf[e4]{Einheit 4 von 4 · Lagebefunde im Sachzusammenhang}
\zweigzeile{keine Prüfungsaufgabe}
\setcounter{aufgabe}{15}

%% TODO Titel: Ich-kann-Satz fehlt in der Bank (Platzhalter: Kettenname) – Nr. 16
%% TODO Anweisung: ein Satz über der ersten Teilaufgabe fehlt in der Bank – Nr. 16
\begin{aufgabe}{Lagebefunde im Sachzusammenhang}
%% TODO \rechenplatz{n}: Schrittzahl des Grundfalls fehlt in der Bank (nur bei fester Schreibform, 2.3 b)
\begin{teile}
\teil Frage: Liegt der Schatten der Lampe auf der Wand? Rechnung: Durchstoßpunkt der Lichtgeraden mit der Wandebene $(2 | 0 | 1{,}8)$. Ist die Frage damit beantwortet? \\ \kreuz{ja} \\ \kreuz{nein – ein Bereich ist noch zu prüfen}
\teil Eine Wand liegt in der Ebene $y = 6$ mit $0 \le x \le 6$ und $0 \le z \le 3$. Eine Lampe in $L(2 | 0 | 4)$ beleuchtet die Spitze $P(3 | 2 | 3)$ eines Pfostens. Liegt der Schatten von $P$ auf der Wand?
\teil Sonnenlicht fällt in Richtung $(-1 | -2 | -1)$; eine Hauswand liegt in der xz-Ebene mit $0 \le x \le 6$ und $0 \le z \le 3$. Schritt I: $(x | 0 | z) = (4 | 2 | 3) + \lambda \cdot (-1 | -2 | -1)$ liefert $\lambda = 1$ und $(3 | 0 | 2)$. Schritt II: $0 < 3 < 6$ und $0 < 2 < 3$. Erläutere beide Schritte im Sachzusammenhang. \hfill (Abitur 2024 GK)
\teil Ein Netz steht in der Ebene $y = 6$, seine Oberkante ist $2$ m hoch (1 LE = 1 m, die xy-Ebene ist der Boden). Ein Ball fliegt von $P(1 | 4 | 3)$ geradlinig in Richtung $(0 | 4 | -1)$. Untersuche rechnerisch, ob er das Netz berührt. \hfill (Abitur 2018 LK)
\teil Ein Spielfeld ist das Rechteck mit $0 \le x \le 9$ und $0 \le y \le 18$ auf dem Boden (xy-Ebene, 1 LE = 1 m). Nach dem Schlag bewegt sich der Ball auf $X_t(4 | 8t + 3 | -5t^2 + 3t + 2)$, $t$ in Sekunden. Untersuche, ob der Ball im Spielfeld auf den Boden trifft. \hfill (Abitur 2018 LK)
\end{teile}
\end{aufgabe}

%% TODO Titel: Ich-kann-Satz fehlt in der Bank (Platzhalter: Kettenname) – Nr. 17
%% TODO Anweisung: ein Satz über der ersten Teilaufgabe fehlt in der Bank – Nr. 17
\begin{aufgabe}{Gerade und Ebene: Länge eines Schattens auf einer Dachfläche über Schnitt von Lichtgerade und Ebene beschreiben}
\begin{teile}
\teil Auf dem schrägen Dach eines Carports, das in der Ebene $D$ liegt, steht im Punkt $F$ eine senkrechte Antenne mit der Spitze $T$. Das Sonnenlicht fällt in parallelen Geraden mit dem Richtungsvektor $\vec{v}$; der Schatten der Antenne liegt ganz auf dem Dach. Beschreibe, wie man die Länge des Schattens berechnet, wenn $T$ und $\vec{v}$ bekannt sind. \hfill (Abitur 2017 LK)
\end{teile}
\end{aufgabe}

%% TODO Titel: Ich-kann-Satz fehlt in der Bank (Platzhalter: Kettenname) – Nr. 18
%% TODO Anweisung: ein Satz über der ersten Teilaufgabe fehlt in der Bank – Nr. 18
\begin{aufgabe}{Lagebefunde im Sachzusammenhang – Fehler finden, Begründen}
\begin{teile}
\teil Ich finde den Fehler: Eine Wand liegt in der Ebene $y = 4$ mit $0 \le x \le 1{,}8$ und $0 \le z \le 2{,}5$. Eine Lampe in $L(0 | 0 | 5)$ beleuchtet $P(1 | 2 | 4)$. Paula berechnet den Durchstoßpunkt $(2 | 4 | 3)$ und schreibt: Der Schatten liegt auf der Wand.
\teil Begründe, warum der Schatten eines Punkts $P$ auf einer Wand der Durchstoßpunkt der Lichtgeraden durch $P$ mit der Wandebene ist.
\end{teile}
\end{aufgabe}

% Abhakseite „Das kann ich“ – Zone nur im Gesamt (\mitzone)
%% TODO Ich-kann-Sätze: Platzhalter wie in den Aufgabendateien
\begin{abhakseite}
\ifdefined\mitzone
\abhakgruppe{Kennst du schon}
\abhak{1}{Ebenengleichungen lesen (Koordinatenform, Normalenvektor ablesen, Parameterform)}
\abhak{2}{Geradengleichungen lesen und den allgemeinen Geradenpunkt bilden (Stützpunkt plus Parameter mal Richtungsvektor)}
\abhak{3}{Skalarprodukt berechnen und Orthogonalität erkennen}
\abhak{4}{Gleichungen lösen}
\abhak{5}{Terme mit einem Parameter umformen und Fallunterscheidungen führen}
\abhak{6}{Punkte im räumlichen Koordinatensystem lesen und Bereiche deuten (Koordinatenschranken, Wandmaße)}
\abhak{7}{Ebenengleichungen lesen (Koordinatenform, Normalenvektor ablesen, Parameterform) – Fehler finden}
\abhak{8}{Ebenengleichungen lesen (Koordinatenform, Normalenvektor ablesen, Parameterform) – selbst rechnen}
\fi
\abhakgruppe{Einheit 1 · Punktprobe und Seitenlage}
\abhak{9}{Wer und wogegen?}
\abhak{10}{Punktprobe und Seitenlage}
\abhak{11}{Punktprobe und Seitenlage – Fehler finden, Begründen}
\abhakgruppe{Einheit 2 · Parameter aus der Lagebedingung}
\abhak{12}{Parameter aus der Lagebedingung}
\abhak{13}{Parameter aus der Lagebedingung – Fehler finden, Begründen}
\abhakgruppe{Einheit 3 · Gerade und Ebene}
\abhak{14}{Gerade und Ebene}
\abhak{15}{Gerade und Ebene – Fehler finden, Begründen}
\abhakgruppe{Einheit 4 · Lagebefunde im Sachzusammenhang}
\abhak{16}{Lagebefunde im Sachzusammenhang}
\abhak{17}{Gerade und Ebene: Länge eines Schattens auf einer Dachfläche über Schnitt von Lichtgerade und Ebene beschreiben}
\abhak{18}{Lagebefunde im Sachzusammenhang – Fehler finden, Begründen}
\end{abhakseite}

\end{document}
